\documentclass[11pt]{letter}

\usepackage{geometry}
\geometry{margin=1in}
\usepackage{setspace}
\usepackage{hyperref}

\signature{Decory Edwards}
\address{Decory Edwards \\ Johns Hopkins University \\ Department of Economics}

\begin{document}

\begin{letter}{Hiring Committee \\ Column \\ San Francisco, CA}

\opening{Dear Hiring Committee,}

I am writing to apply for the New Business position for 2026 graduates at Column. I am currently a PhD candidate in Economics at Johns Hopkins University. My training combines rigorous analytical reasoning, structured research, and clear communication. I am motivated by work that requires understanding complex businesses quickly, identifying high potential opportunities, and engaging directly with decision makers. Column’s approach to building a modern banking and software platform, and the opportunity to help shape its go to market motion from the ground up, strongly align with my interests and skill set.

My research agenda is at the intersection of wealth inequality and macroeconomics. In my job market paper, I build and estimate a structural heterogeneous agent model with returns that differ across households. The core contribution is to show how heterogeneity in returns and income risk jointly shape the wealth distribution and consumption behavior. Relative to closely related work, my framework performs better at matching the lower and middle portions of the wealth distribution. This difference matters quantitatively. Models that focus primarily on returns heterogeneity can fit the top of the wealth distribution closely but tend to generate too much wealth among the bottom half of households. In my framework, income uncertainty generates a precautionary saving motive that produces genuinely low wealth households. This leads to a realistic distribution of marginal propensities to consume and aggregate consumption responses that align with empirical estimates.

A key feature of my research is sourcing and synthesizing information from incomplete and noisy environments. I regularly dig into new domains, map industries, evaluate business models, and develop informed views based on limited initial information. I am comfortable conducting outbound research, tailoring messaging to specific audiences, and using conversations to refine my understanding of fit and opportunity. This work mirrors the core responsibilities of new business development and top of funnel qualification.

I am particularly interested in Column because of its focus on financial infrastructure and APIs and its ambition to simplify the banking supply chain for builders and developers. The opportunity to engage directly with founders and executives across a wide range of companies, from early stage startups to public firms, is especially appealing. I value environments that reward curiosity, ownership, and precision, and where quality of insight matters more than volume of outreach.

My economics training has provided a strong foundation in analytical thinking, independent problem solving, and clear written and verbal communication. I am comfortable working with ambiguity, moving quickly, and collaborating closely with sales and partnerships teams to advance qualified opportunities. I take pride in being resourceful, coachable, and focused on team success.

I am enthusiastic about the opportunity to work on site in San Francisco and to contribute in person to Column’s collaborative culture. I see this role as an opportunity to apply my analytical background to business development in a high growth fintech environment while learning directly from experienced operators.

I would welcome the opportunity to contribute my skills and perspective to Column. Thank you for your time and consideration.

\closing{Sincerely,}

\end{letter}
\end{document}
